\documentclass[11pt]{article}
\usepackage[margin=1in]{geometry}
\usepackage{amsmath}
\usepackage{xcolor}
\usepackage{parskip}
\usepackage{hyperref}
\usepackage{microtype}

\newenvironment{reviewer}{\color{blue!50!black}\itshape}{\normalcolor}
\newcommand{\response}{\textbf{Response:} }

\begin{document}

\begin{center}
{\Large \textbf{Response to Reviewers}}\\[6pt]
Manuscript jrfm-4438773\\
\emph{Selective Prediction and the Persistence Illusion: A Diagnostic Decomposition of VIX
Regime Classification}\\[4pt]
Akshat Gupta and Jianguo Liu
\end{center}

We thank the editor and both reviewers for exceptionally constructive reports. The revision
implements every requested experiment and reframes the paper around the diagnostic protocol as
its contribution. We respond point by point below; all section and table numbers refer to the
revised manuscript. We begin with the editor's first major comment, since it is the one point
that could have affected existing numbers---and, in the interest of full transparency, the
editor's concern was warranted.

\section*{Response to the Editor}

\subsection*{Major Comments}

\textbf{Major 1 (calibration of the persistence abstention constant $c$).}
\begin{reviewer}
If $c$ is calibrated on the test set, the baseline is tuned ex post on evaluation data\ldots
Please state explicitly how $c$ is chosen, and if it uses test-set coverage, re-run with $c$
set on the training window.
\end{reviewer}

\response We audited the code and report the finding plainly: in the submitted version, $c$
(and the analogous coverage-matching thresholds for the LR, HAR-LR, and Markov baselines) was
selected by matching the RF's \emph{test-window} covered count, using test-window VIX levels
and model probabilities. No test \emph{labels} entered any threshold choice---the matching
targeted coverage counts, not accuracy---but the editor is correct that this constituted
ex-post use of evaluation data, and we agree it was inappropriate in a paper about evaluation
artifacts. We have recalibrated every baseline threshold on the training window only: each
threshold is now the training-window quantile of the relevant confidence measure that
reproduces the RF's \emph{training} coverage rate at $\tau = 0.25$, frozen before the test
window is touched. This procedure is now stated explicitly in Section 5.2 and Appendix F. We
re-ran every affected analysis (baseline comparison, McNemar, bootstrap, 2022-exclusion,
S\&P 500 replication, walk-forward folds), and additionally fixed the same pattern in two
places the review could not have seen: the bootstrap helper function and the per-fold
walk-forward persistence comparison. A full leakage sweep confirmed the remaining pipeline is
clean (scalers fit on training only; calibration via \texttt{TimeSeriesSplit} within training;
backward-looking feature windows; hyperparameters fixed a priori or tuned on training only).

The changes to the numbers are reported honestly in the revision. The McNemar result is
essentially unchanged ($b = c = 0$ at $N \in \{10, 20, 25\}$; disagreement on at most 3 of 708
jointly covered days elsewhere, always in the RF's favor). Calm$\to$high covered accuracy
remains 0\% everywhere. Two results changed materially: (i) the RF-minus-persistence gap at
$N = 10$ is now statistically significant ($+2.8$ pp, 95\% CI $[+0.18, +5.27]$), which the
revised Section 5.9 attributes---via the unchanged McNemar result---to coverage composition
rather than predictive disagreement; and (ii) the HAR-LR no longer exceeds the RF in point
estimate at four horizons (only at $N = 25$), so that claim has been replaced throughout by
the statistically supported statement that the HAR-minus-RF bootstrap CIs include zero. We
believe the corrected comparison is more honest and, because the coverage-composition
mechanism is now demonstrated rather than asserted, ultimately more persuasive.

\textbf{Major 2 (coverage matching $\neq$ composition matching; McNemar context).}

\response Fully agreed. New Section 5.10 (``Composition of the Jointly Covered Set'') and
Table 12 report, for each horizon, the size of RF $\cap$ persistence as a fraction of test
days and its breakdown across the four transition types. The joint set is dominated by
regime-stable days (92.9\% at $N = 5$; 92.9\% calm$\to$calm alone at $N = 25$). The revised
text states the implication in the editor's own spirit: the $b = c = 0$ result is close to
guaranteed \emph{given} that composition---the filter selects easy days, and that is the
finding, not a nuisance.

\textbf{Major 3 (circularity between label design and the persistence benchmark).}

\response New Section 5.3 adds the requested continuous-forecast-then-threshold benchmark: a
HAR regression of $\text{VIX}_{t+N}$ on daily/weekly/monthly VIX components, estimated on the
training window, thresholded at 20, with an abstention band calibrated on training forecasts
and frozen. Its behavior reproduces the discrete pipeline in full---overall covered accuracy
within 1--7 pp of the RF and 0\% covered accuracy on calm$\to$high transitions at every
horizon (Table 5). The diagnostic structure is therefore not an artifact of binarizing a
mean-reverting process, and the revised text says exactly this.

\textbf{Major 4 (bootstrap block length).}

\response We re-ran the moving-block bootstrap at block lengths 32, 60, and 90 days
(Appendix D, Table 22). The intervals are qualitatively stable: the $N = 10$ interval excludes
zero at all three block lengths, and every other interval includes zero at all three. Section
5.9 now acknowledges the dependence-order argument explicitly.

\textbf{Major 5 (single-threshold evaluation; risk--coverage curves).}

\response New Section 5.11 reports full risk--coverage curves per horizon for RF, HGB,
XGBoost, and the persistence rule, with AURC in Table 13 (Figure 8). The threshold-free view
adds an honest nuance, which we report: the ML models rank their own confidence better than
the distance-to-threshold heuristic over the full coverage range (RF attains the lowest AURC
at four of five horizons), but this ranking advantage operates within the regime-stable
population and does not surface transition days at any coverage level.

\textbf{Major 6 (effective sample on the key result; episodes).}

\response New analysis in Section 5.5 clusters covered calm$\to$high days into episodes
(gap $>$ 10 trading days). The counts are 16 episodes at $N = 5$, 15 at $N = 10$, 12 at
$N = 15$, and 4--5 at $N = 20$ and 25, with dates listed. At short horizons the 0\% record
therefore spans substantially more independent events than the 4--6 the decision letter
anticipated; at the two longest horizons we temper the language exactly as requested (the
revised text reads ``consistent with low---but not precisely estimated---transition skill'').

\textbf{Major 7 (AUC and Brier by horizon).}

\response Added as Appendix E (Table 23): AUC and Brier, pre- and post-calibration, for RF,
HGB, and XGBoost at all horizons, with an interpretive paragraph noting that aggregate AUC on
a 70/30 test set is itself dominated by the regime-stable population.

\subsection*{Minor Comments}

\textbf{Minor 1 (selective-prediction evaluation literature).} Added to Section 2.3: El-Yaniv
and Wiener (2010) on risk--coverage foundations, Geifman et al.\ (2019) on AURC, and Fisch et
al.\ (2022) on selective calibration, with a sentence relating our protocol to that strand.

\textbf{Minor 2 (unused XGBoost reference).} Now cited in Section 4.2, where XGBoost joins the
model set.

\textbf{Minor 3 (per-horizon sample sizes).} Table 1 now reports train/test $N$ and class
frequencies for every horizon.

\textbf{Minor 4 (``informational ceiling'').} Retired as a global claim. The abstract and
conclusions now use the Section 5.7 register: the failure is informational \emph{with respect
to this feature set}, and richer information sets are flagged as the open question.

\textbf{Minor 5 (AI declaration).} We appreciate the comment; the declaration has been updated
to reflect the actual scope of assistance (analysis-code development, data processing, and
manuscript drafting), with the authors' verification of all results stated explicitly.

\textbf{Code release.} As encouraged, the complete pipeline is now public at
\url{https://github.com/akshatg10104/VIX_research} and cited in the Data Availability
Statement.

\section*{Response to Reviewer 1}

\textbf{Major 1 (novelty positioning).}

\response The Introduction and Conclusions have been rewritten to lead with the diagnostic
protocol as the contribution, with the persistence result as its demonstration case. A closing
Introduction paragraph now separates (i) what this study demonstrates (two highly persistent
financial classification tasks) from (ii) the broader hypothesis about selective prediction
under serial dependence, for which we claim only two supporting cases.

\textbf{Major 2 (GARCH / regime-switching benchmarks).}

\response Two changes. First, Section 5.3 adds a continuous-forecast-then-threshold HAR
benchmark, which empirically represents the autoregressive forecasting class that Fernandes et
al.\ (2014) show is competitive for the VIX level; it reproduces the transition failure
exactly. Second, Section 5.2 adds an explicit justification paragraph: GARCH-family models
target the conditional variance of returns rather than the option-implied index level, and a
GARCH-to-VIX-threshold conversion would require variance-risk-premium assumptions that make
the comparison ambiguous, while Markov-switching dynamics are already represented by the
Markov-chain baseline.

\textbf{Major 3 (methodological transparency).}

\response The complete code is now public (see above), and new Appendix F documents each item
raised: feature construction timing, missing-observation handling, hyperparameter selection,
absence of indirect test-set tuning, per-horizon calibration, baseline threshold calibration,
and the exact bootstrap implementation.

\textbf{Major 4 (statistical inference; confirmatory vs.\ descriptive).}

\response The Conclusions now open with a two-tier classification: statistically supported
findings (McNemar; bootstrap CIs; the 0\% transition record and its episode count) versus
descriptive observations (point-estimate orderings within overlapping CIs), with an explicit
multiple-comparisons acknowledgment noting that the central claims are negative results, for
which multiplicity cuts against us finding spurious \emph{significance} rather than in favor
of our conclusions.

\textbf{Major 5 (generalizability).}

\response All ``structural property'' language has been rescoped to the two studied problems;
the general claim is now stated as a hypothesis, and cross-domain replication (credit spreads,
recession indicators, currency crises) is listed in Future Work as required before broader
claims.

\textbf{Major 6 (practical implications).}

\response New subsection ``Practical Implications: Who Should Use This Protocol, and How'' in
the Conclusions names four audiences (model validators/MRM teams, quantitative researchers,
referees and editors, practitioners), gives a six-item reporting checklist derived from the
protocol, and discusses when abstention still adds value---including the abstention-spike
leading-indicator observation as a worked example.

\textbf{Minor comments 1--8.}

\response All implemented: (1) the Introduction now explicitly distinguishes forecasting
volatility from forecasting binary volatility regimes; (2) Section 3.3 expands the
threshold-20 motivation (Whaley convention, practitioner usage), points to the 18/22
robustness table, and notes data-driven threshold selection as an alternative; (3) per-horizon
class frequencies appear in Table 1; (4) the threshold-sensitivity, VIX-threshold-robustness,
and sustained-label tables have been moved to Appendix B to streamline the main text; (5)
figure and table captions have been audited to state the metric and its base population
(covered vs.\ all days) throughout; (6) Section 5.8 now notes that SHAP measures model
attribution, not causal importance; (7) Section 5.7 adds the LSTM architecture justification
(parameter budget matched to the ensembles, window matched to the longest feature lookback)
and notes that architectures were deliberately not swept; (8) the global wording pass replaced
``confirms/proves/establishes'' with ``suggests/indicates/is consistent with'' except where a
result is a deterministic identity.

\section*{Response to Reviewer 2}

\textbf{Comment 1 (abstract too long).}

\response The abstract has been shortened by roughly a third and now leads with the protocol,
retaining three headline results.

\textbf{Comment 2 (persistence dominating ML is not surprising; need a methodological
contribution).}

\response We agree the phenomenon is known, and the revision no longer presents it as the
contribution. The contribution, stated in the first page of the Introduction, is the reusable
six-component diagnostic protocol---coverage-matched baselines calibrated on training data,
joint-coverage composition analysis, transition-conditional auditing, risk--coverage
evaluation, attribution analysis, and abstention confound testing---together with its
demonstration on two markets. To our knowledge no comparable protocol exists for selective
prediction on serially correlated labels.

\textbf{Comment 3 (stronger ML baselines: XGBoost, LightGBM, CatBoost, Transformers, TFT).}

\response Following the editor's guidance on scope (``I would therefore accept either (a)
adding one strong boosted model plus a continuous-forecast-then-threshold comparison\ldots''),
we added XGBoost through the identical pipeline (Section 5.3 and Appendix C) plus the HAR
forecast-then-threshold benchmark. XGBoost's selective accuracy, coverage, and
balanced-accuracy profile track the RF closely---and its calm$\to$high covered accuracy is
0.0\% at all five horizons, exactly as the paper's informational-constraint logic predicts for
any architecture trained on this feature set. We believe this demonstration makes additional
boosted variants uninformative, per the editor's ruling.

\textbf{Comment 4 (richer predictors: option-implied measures, news, order flow).}

\response Two clarifications. First, the feature set already contains option-implied
measures: the CBOE SKEW index (three features) and the VIX term-structure spread
(VIX3M $-$ VIX) are both derived from index option prices, so the negative result already
covers the most accessible option-implied information. Second, the restriction to standard
daily public data is deliberate: it makes the negative result a clean statement about a
well-defined information set. The Limitations and Future Work sections now state prominently
that news sentiment, order flow, intraday options flow, and macro surprise indices are the
primary open direction for testing whether transitions are predictable from richer data.

\textbf{Comment 5 (robustness: more out-of-sample testing, other markets/periods).}

\response We respectfully note that the submitted manuscript already contained several of the
requested checks, which the revision retains and extends: 5-fold expanding walk-forward
validation (Section 5.6); a full cross-market replication on S\&P 500 trend regimes
(Section 6); VIX-threshold robustness at 18/20/22 (Appendix B); and a 2022-bear-market
exclusion check (Section 5.9). The revision adds: bootstrap block-length sensitivity at
32/60/90 days (Appendix D), threshold-free risk--coverage curves with AURC (Section 5.11),
and the XGBoost and HAR forecast-then-threshold benchmarks (Section 5.3). Extension to
non-U.S.\ volatility indices (e.g., VSTOXX) is noted in the Limitations as future work.

\textbf{Comment 6 (practical implications; who can use these results).}

\response Addressed by the new Practical Implications subsection described in our response to
Reviewer 1, Major 6: four named audiences and a concrete six-item evaluation checklist.

\vspace{1em}
We are grateful for the care taken by the editor and both reviewers; the revision is
substantially stronger for it. We hope the revised manuscript now meets the standards of the
journal.

\begin{flushleft}
Sincerely,\\[4pt]
Akshat Gupta and Jianguo Liu
\end{flushleft}

\end{document}
